\SGARefTarget{sga3:xviii:section:2}\section*{2. Local Determination of a
Group Homomorphism}
\addcontentsline{toc}{section}{2. Local Determination of a Group
Homomorphism}

Let \(G\) and \(H\) be groups, and let \(U\subseteq G\) be a subset
such that \(U\cdot U=G\). If \(f\) and \(g\) are homomorphisms from
\(G\) to \(H\) such that \(f=g\) on \(U\), then \(f=g\). Likewise:

\SGARefTarget{sga3:xviii:prop:2.1}\paragraph{\SGARefLink{sga3:xviii:prop:2.1}{Proposition 2.1}.}
\textit{Let \(S\) be a site (cf.\ Exposé IV),
let \(G\) be a sheaf of groups on \(S\), and let \(U\subseteq G\) be
a subsheaf of sets such that the morphism \(U\times U\to G\)
induced by multiplication is an epimorphism. If \(f\) and \(g\) are
homomorphisms from \(G\) to a sheaf of groups \(H\) and are equal on
\(U\), then \(f=g\).}

\SGARefTarget{sga3:xviii:cor:2.2}\paragraph{\SGARefLink{sga3:xviii:cor:2.2}{Corollary 2.2}.}
\textit{Let \(G\) be an \(S\)-group prescheme that is locally of
finite presentation and flat, let \(U\subseteq G\) be a relatively
schematically dense open subset, and let \(H\) be a sheaf of groups
on \(S\) for the fppf topology. Then every homomorphism
\(f:G\to H\) is determined by its restriction to \(U\).}

Indeed, since \(G\) is flat over \(S\), the composition law on \(G\)
is a flat morphism (VI\textsubscript{B}, \SGARefLink{concl:sga3-VIB-9.2-xi}{9.2(xi)}); it follows that
\(U\times_S U\to G\) is faithfully flat and locally of finite
presentation, hence an fppf covering, and therefore an epimorphism.

\SGARefTarget{sga3:xviii:prop:2.3}\paragraph{\SGARefLink{sga3:xviii:prop:2.3}{Proposition 2.3}.}
\textit{Let \(G\) be a group prescheme locally of finite
presentation and flat over \(S\), let \(U\subseteq G\) be a
relatively schematically dense open subset, and let \(H\) be a
sheaf of groups for the fppf topology. Write \(m_G\) for the
multiplication on \(G\) and \(m_H\) for that on \(H\).}

\textit{Let \(\bar f:U\to H\) be an \(S\)-morphism, and suppose that
there exists a relatively schematically dense open subset \(V\) of
\(m_G^{-1}(U)\cap(U\times_S U)\) such that the following diagram is
commutative:}

\begin{center}
\begin{tikzcd}[column sep=large,row sep=large]
  V \arrow[r,"{(\bar f\times\bar f)|_V}"] \arrow[d,"m_G"] &
  H\times_S H \arrow[d,"m_H"]\\
  U \arrow[r,"\bar f"] & H
\end{tikzcd}

\end{center}

\textit{Then there exists a morphism of \(S\)-groups
\(f:G\to H\), necessarily unique, extending \(\bar f\), in each of
the following situations:}

\begin{enumerate}
\item[\SGAInlineRefTarget{sga3:xviii:prop:2.3:item:i}(i)] \textit{\(H\) is representable.}
\item[\SGAInlineRefTarget{sga3:xviii:prop:2.3:item:ii}(ii)] \textit{The diagonal morphism
  \(H\to H\times_S H\) is representable by a closed
  immersion.}
\item[\SGAInlineRefTarget{sga3:xviii:prop:2.3:item:iii}(iii)] \textit{For every section \(a\in U(S)\), the open
  subset
  \(\operatorname{pr}_2((a\times_S U)\cap V)\) is relatively
  schematically dense in \(U\)\begingroup
  \renewcommand{\thefootnote}{*}\footnote{Or in \(G\), which is
  equivalent by \SGARefLink{sga3:xviii:prop:1.1:item:iv}{1.1(iv)}.}
  \addtocounter{footnote}{-1}
  \endgroup, and this assertion remains true after every base
  change \(S'\to S\).}
\end{enumerate}

\begin{proof}
First note that \(V\) is relatively schematically dense in
\(G\times_S G\). Indeed, \(V\) is relatively schematically dense in
\[
  m_G^{-1}(U)\cap(U\times_S G)\cap(G\times_S U),
\]
and each of the three factors is obtained from \(U\subseteq G\) by
an evident base change \(G\times_S G\to G\). The assertion therefore
follows from \SGARefLink{sga3:xviii:prop:1.1}{1.1}.

To construct \(f:G\to H\), since \(U\times_S U\to G\) is an fppf
covering, it is enough to construct a morphism \(U\times_S U\to H\)
whose two pullbacks to
\[
  (U\times_S U)\times_G(U\times_S U)
\]
are equal. Take the morphism
\[
  m_H\circ(\bar f\times\bar f):U\times_S U\longrightarrow H.
\]
For sections \(a,b,c,d\in U(S')\), where \(S'\to S\) is arbitrary,
such that \(ab=cd\), we must prove
\[
  \bar f(a)\bar f(b)=\bar f(c)\bar f(d).
\]
If \((a,b)\) and \((c,d)\) belong to \(V(S')\)\begingroup
\renewcommand{\thefootnote}{\ensuremath{\ddagger}}\footnote{The
French source prints \((a,d)\) and \((c,d)\). The first pair must be
\((a,b)\), as required by the displayed equality and by the
compatibility hypothesis on \(V\).}
\addtocounter{footnote}{-1}
\endgroup, the desired equality follows at once. Thus it holds on
\((V\times_G V)(S')\).

I claim that \(V\times_G V\) is a relatively schematically dense
open subset of
\[
  (U\times_S U)\times_G(U\times_S U).
\]
This will complete the proof in cases \textup{\SGARefLink{sga3:xviii:prop:2.3:item:i}{(i)}} and \textup{\SGARefLink{sga3:xviii:prop:2.3:item:ii}{(ii)}}
by \SGARefLink{sga3:xviii:prop:1.4}{1.4}, because the morphism \(f\) constructed in this way extends
\(\bar f\), and its being a homomorphism is then evident, again by
\SGARefLink{sga3:xviii:prop:1.4}{1.4}. Now
\[
  V\times_G V
  =
  \bigl(V\times_G(G\times_S G)\bigr)
  \cap
  \bigl((G\times_S G)\times_G V\bigr).
\]
By symmetry and \SGARefLink{sga3:xviii:prop:1.1}{1.1}, it is enough to check that
\(V\times_G(G\times_S G)\) is relatively schematically dense in
\((G\times_S G)\times_G(G\times_S G)\). The latter prescheme is
isomorphic over \(S\) to \(G\times_S G\times_S G\), the isomorphism
being given by
\[
  (a,b,c,d)\longmapsto(a,b,c).
\]
What remains is therefore the relative schematic density of
\(V\times_S G\) in \(G\times_S G\times_S G\), which follows from
that of \(V\) in \(G\times_S G\).

It remains to treat case \textup{\SGARefLink{sga3:xviii:prop:2.3:item:iii}{(iii)}}. To prove
\(\bar f(a)\bar f(b)=\bar f(c)\bar f(d)\), we may make an fppf
covering base change. Suppose that, after such a base change
\(S'\to S\), there is a section \(x\in G(S')\) such that
\[
  (b,x),\quad(a,bx),\quad(d,x),\quad(c,dx)
\]
all belong to \(V\). Then
\[
\begin{aligned}
  \bar f(a)\bar f(b)\bar f(x)
    &=\bar f(a)\bar f(bx)
     =\bar f(abx)
     =\bar f(cdx)\\
    &=\bar f(c)\bar f(dx)
     =\bar f(c)\bar f(d)\bar f(x),
\end{aligned}
\]
which gives the desired equality.

To find such an \(x\)\footnote{Editor's note: The beginning of the
following sentence has been added.}, set, for every \(z\in U(S')\),
\[
  V_z=
  \operatorname{pr}_{2/S'}
  \bigl((z\times_{S'}U_{S'})\cap V_{S'}\bigr).
\]
The conditions imposed on \(x\) say exactly that \(x\in W\), where
\[
  W=
  V_b(S')\cap V_d(S')\cap
  b^{-1}V_a(S')\cap d^{-1}V_c(S').
\]
By hypothesis \textup{\SGARefLink{sga3:xviii:prop:2.3:item:iii}{(iii)}}, \(W\) is relatively schematically
dense in \(G\). \SGARefLink{sga3:xviii:prop:1.7}{Proposition 1.7} therefore supplies such a section
\(x\) after an fppf covering base change.

Let us verify that the morphism thus constructed is multiplicative.
Let \(a,b\in G(S')\), where \(S'\to S\) is arbitrary; after this
base change, rename \(S'\) as \(S\). Choose an fppf covering base
extension \(S'\to S\) and a section \(x\in G(S')\) such that
\[
  x\in U(S')
  \qquad\text{and}\qquad
  ax^{-1}\in U(S'),
\]
that is, \(x\in(a^{-1}U)^{-1}(S')\). Choose in addition another
fppf covering extension \(S''\to S'\) and a section
\(y\in G(S'')\) such that
\[
  y,\quad by^{-1},\quad aby^{-1}\in U(S''),
\]
and
\[
  by^{-1}\in V_x,
  \qquad
  xby^{-1}\in V_{ax^{-1}}.
\]
Then, by the definition of \(f\), over \(S''\) we have
\[
  f(a)=\bar f(ax^{-1})\bar f(x),\qquad
  f(b)=\bar f(by^{-1})\bar f(y),\qquad
  f(ab)=\bar f(aby^{-1})\bar f(y).
\]
Furthermore,
\[
\begin{aligned}
  f(a)f(b)
    &=\bar f(ax^{-1})\bar f(x)\bar f(by^{-1})\bar f(y)\\
    &=\bar f(ax^{-1})\bar f(xby^{-1})\bar f(y)\\
    &=\bar f(aby^{-1})\bar f(y)
     =f(ab),
\end{aligned}
\]
which proves multiplicativity.

It is now easy to see that \(f\) extends \(\bar f\). Let
\(a\in U(S')\); after base change, rename \(S'\) as \(S\), and choose
an fppf covering \(S'\to S\) and sections \(x,y\in G(S')\) such that
\[
  (x,y),\qquad (ax,y),\qquad (a,xy)
\]
belong to \(V(S')\). Then
\[
  f(xy)=\bar f(x)\bar f(y)=\bar f(xy),
  \qquad
  f(axy)=\bar f(ax)\bar f(y)=\bar f(axy).
\]
Consequently,
\[
\begin{aligned}
  f(a)
    &=f(axy)f\bigl((xy)^{-1}\bigr)
     =f(axy)f(xy)^{-1}\\
    &=\bar f(axy)\bar f(xy)^{-1}
     =\bar f(a).
\end{aligned}
\]
\end{proof}

\SGARefTarget{sga3:xviii:remark:2.3.1}\paragraph{\SGARefLink{sga3:xviii:remark:2.3.1}{Remark 2.3.1}.}
In many cases hypothesis \textup{\SGARefLink{sga3:xviii:prop:2.3:item:iii}{(iii)}} holds, for it will still
hold if \(U\) is replaced by a smaller open subset \(U'\), still
relatively schematically dense in \(G\). For example:

\SGARefTarget{sga3:xviii:prop:2.4}\paragraph{\SGARefLink{sga3:xviii:prop:2.4}{Proposition 2.4}.}
\textit{In the situation of \SGARefLink{sga3:xviii:prop:2.3}{2.3}, suppose that every geometric fiber
of \(G/S\) is irreducible. Put
\[
  U'=\operatorname{pr}_1V
  \qquad\text{and}\qquad
  V'=V\cap m_G^{-1}(U')\cap(U'\times_S U').
\]
Then \(U'\subseteq U\) is an open subset relatively schematically
dense in \(G\), and \(U'\), \(\bar f|_{U'}\), and \(V'\) satisfy
hypothesis \textup{\SGARefLink{sga3:xviii:prop:2.3:item:iii}{(iii)}}.}

\begin{proof}
The subset \(U'\) is open because
\(G\times_S G\to G\) is flat and locally of finite presentation.
All the other verifications are immediate except for hypothesis
\textup{\SGARefLink{sga3:xviii:prop:2.3:item:iii}{(iii)}}.

Let \(a\in U(S')\); after base change, rename \(S'\) as \(S\). To
verify that
\[
  \operatorname{pr}_2\bigl((a\times_S U')\cap V'\bigr)
\]
is relatively schematically dense in \(U'\), it suffices by
\SGARefLink{sga3:xviii:cor:1.3}{Corollary 1.3} to do so fiber by fiber. Thus it is enough to treat the
case where \(S\) is the spectrum of a field, and in that case it is
enough to prove that \(U'\) is nonempty, since \(G\) is irreducible
and has no embedded components (cf.\ VI\textsubscript{A}, \SGARefLink{prop:sga3-VIA-1.1.1}{1.1.1}).
Now
\[
\begin{aligned}
  \operatorname{pr}_2\bigl((a\times_S U')\cap V'\bigr)
  &=
  \operatorname{pr}_2\bigl((a\times_S U')\cap V\bigr)\\
  &\quad{}\cap
  \operatorname{pr}_2\bigl((a\times_S U')\cap m_G^{-1}(U')\bigr).
\end{aligned}
\]
The second factor on the right is dense in \(G\). It therefore
remains only to prove that
\[
  \operatorname{pr}_2\bigl((a\times_S U')\cap V\bigr)
\]
is dense in \(G\), that is, nonempty; and this is clear because
\(a\in U'=\operatorname{pr}_1V\).
\end{proof}
